\section{Mathematical Analysis}
\label{app:math}
% Optional compact mathematical appendix. Add at most what helps the story;
% the main text already contains the interaction and gate cost identities.
\subsection{Agreement Bound}
For nonempty binary masks viewed as finite pixel sets, let $J$ be IoU and
$d(X,Y)=1-J(X,Y)$ be Jaccard distance~\cite{kosub2019}.
The reverse triangle inequality gives
\begin{equation}
  \left|J(M_t,M_q)-J(M_v,M_q)\right|
  \leq 1-J(M_t,M_v).
  \label{eq:agreement-bound}
\end{equation}
Indeed, applying the triangle inequality in both directions bounds
$|d(M_t,M_q)-d(M_v,M_q)|$ by $d(M_t,M_v)$.
High agreement therefore bounds the difference between the two probe
accuracies, rather than their absolute accuracy. For example, if both
probes predict the same region disjoint from the target, their mutual
IoU is one while both target IoUs are zero. The bound concerns the two
probe masks, not the subsequent joint-prompt prediction.

\subsection{IoU Aggregation}
\balance
Within one class, let $U_e=\mathrm{TP}_e+\mathrm{FP}_e+\mathrm{FN}_e$
be episode union size and $J_e=\mathrm{TP}_e/U_e$. With population
covariance over episodes,
\begin{equation}
  \begin{aligned}
  J_{\mathrm{pool}}-\overline{J}
  &=\frac{\operatorname{Cov}(U,J)}{\overline{U}}\\
  &=\frac{\operatorname{Cov}(G,J)+\operatorname{Cov}(\mathrm{FP},J)}
         {\overline{U}},
  \end{aligned}
  \label{eq:pooling-covariance}
\end{equation}
where $G_e=\mathrm{TP}_e+\mathrm{FN}_e$ is target area.
The identity follows from $\mathrm{TP}_e=U_eJ_e$ and
$\overline{UJ}=\overline{U}\,\overline{J}+\operatorname{Cov}(U,J)$.
It explains why a large erroneous foreground region can lower pooled
IoU much more than the mean episode IoU: the error also increases its
weight through $U_e$.

For CG-ICS on LVIS, the mean episode IoU is $55.19\%$ and the class-pooled
IoU is $47.44\%$. Averaging Eq.~\ref{eq:pooling-covariance} using the
same class/fold weights yields a $-7.75$-point gap: the target-area
covariance term is $+3.82$ points and the false-positive covariance
term is $-11.57$ points. These are exact weighting components, not
estimates of causal effects.


